The statement as written in the prompt is ambiguous: if we literally define ``$v(k)$ is the minimal integer that does \emph{not} appear as a denominator in any $k$--term Egyptian fraction decomposition of $1$,'' then $1$ never appears as a denominator when $k\ge 2$ (because $1=1/1$ uses only one term), hence $v(k)=1$ for all $k\ge2$, which is trivial and not the intended problem.

Following the standard formulation used in the literature on this Erd\H{o}s problem, I will interpret $v(k)$ as the \emph{smallest integer $>1$} that does not appear.

\medskip

\noindent\textbf{Definitions.}
\begin{itemize}
\item For $k\ge 1$, let $\mathcal{D}_k$ be the set of integers $n>1$ for which there exists a strictly increasing $k$--tuple of integers
\[
2\le n_1< n_2<\cdots<n_k,\qquad \sum_{i=1}^k \frac1{n_i}=1,
\]
with $n\in\{n_1,\dots,n_k\}$.
\item Define
\[
 v(k) := \min\{n>1: n\not\in \mathcal{D}_k\}.
\]
\end{itemize}

\noindent\textbf{Goal.} Estimate the growth of $v(k)$ as $k\to\infty$ (e.g. find the best possible lower/upper bounds in terms of $k$, and ideally its correct order of magnitude).
